\begin{table}[t]
\centering
\caption{Accounting for the street gap.}
\label{tab:gap_accounting}
\scriptsize
\begin{threeparttable}
\begin{tabular}{@{\extracolsep{0pt}}lrrrrrr}
\toprule
 & \multicolumn{2}{c}{Sex ratio (F/1000 M)} & & \multicolumn{3}{c}{Share of the gap} \\
\cmidrule(lr){2-3}\cmidrule(lr){5-7}
City & Residential & Observed & Shortfall & Not out & Fewer trips & Residual \\
\midrule
Mumbai & 853 & 236 & 72\% & 62\% & 14\% & 24\% \\
Navi Mumbai & 837 & 221 & 74\% & 61\% & 14\% & 25\% \\
Bangalore & 923 & 377 & 59\% & 76\% & 17\% & 7\% \\
Delhi & 876 & 246 & 72\% & 63\% & 14\% & 24\% \\
\bottomrule
\end{tabular}
\begin{tablenotes}[flushleft]
\item Shortfall is the headline comparison of the observed pedestrian sex ratio to the residential sex ratio. It is reported unadjusted.
\item The three shares decompose that shortfall. Women are less likely to leave home at all (47.3\% against 86\% of men) and, having left, make fewer trips (1.32 against 2.93 per day; 0.82 conditional on being mobile). Both rates are for urban India from the 2019 Time Use Survey \citep{goel2023gendergap}. Scaling the residential ratio by each factor in turn partitions the gap.
\item The partition locates the gap; it does not explain it. The headline is reported unadjusted because women kept at home are the largest component of the absence, not a confound to remove.
\item The residual is the part only street imagery can see: women who left home, made their trips, and are still not among pedestrian sightings. A time-use survey records who left the house, not who is standing on a road.
\item The mobility rates are national urban averages applied to four specific cities. If women in these cities are more mobile than the urban average, the front-door and trip-count shares are overstated and the residual is a lower bound.
\item Scaling a residential ratio by $\pm$10\% moves the residual share by at most 8 percentage points.
\end{tablenotes}
\end{threeparttable}
\end{table}
